\chapter{Калибровочная инвариантность и необходимость векторного потенциала}

Прежде чем переходить к физике, нужно обсудить 2 практических вопроса. Следующее занятие будет пропущено, пока лектор находится в Бельгии, но курс продолжится, пока не будут прочитаны все 10 лекций. Рукописные конспекты к сегодняшнему вечеру тоже были подготовлены в спешке; если кто-нибудь согласится отсканировать их для сайта курса, это спасёт их от поразительно медленного сканера.

Предыдущая лекция была математически насыщенной. Вместо того чтобы делать вид, будто это не так, мы ещё раз пройдём её центральное вычисление. Цель курса — подлинный теоретический минимум: достаточно структуры, чтобы перейти на следующий уровень. Если бы математику можно было убрать без потери физики, были бы все основания её убрать. Но точные законы требуют количественных утверждений, а количественные утверждения требуют математики. Обозначения могут стать мудрёными; это не свидетельствует о том, что физики пытаются скрыть предмет за жреческим языком.

Возможно, ещё предстоит найти более простую формулировку законов природы. Можно, как шутил Фейнман, возвести в квадрат каждое уравнение физики, сложить квадраты и приравнять сумму к нулю. Тогда получится одно уравнение, потому что каждый квадрат должен будет обратиться в нуль, но нового понимания это не даёт. Простота означает нахождение организующей идеи, а не просто сжатие записи. Когда кто-то из аудитории возражает против предположения, что такая идея непременно должна прийти от кого-то молодого, с этим легко соглашаются: возраст не имеет отношения к проницательности.

\lectureoverview{Сначала мы восстановим индексные обозначения, необходимые для распознавания лоренцевых скаляров. Затем сформулируем 4 принципа релятивистской теории, полностью повторим вывод для заряженной частицы и, наконец, дважды обнаружим калибровочную инвариантность: как граничный член в действии и как инвариантность тензора электромагнитного поля.}

\section{Возвращение к индексному языку}

\subsection{Почему некоторые индексы пишутся сверху}

Пусть
\begin{equation}
x^\mu=(t,x,y,z),
\qquad
\mu=0,1,2,3,
\end{equation}
с метрикой
\begin{equation}
\eta_{\mu\nu}
=
\operatorname{diag}(-1,1,1,1).
\label{eq:sr7-metric}
\end{equation}
Координатам и координатным смещениям приписываются верхние индексы:
\begin{equation}
dx^\mu=(dt,dx,dy,dz).
\end{equation}
Производные по этим координатам несут нижние индексы:
\begin{equation}
\partial_\mu
\equiv
\frac{\partial}{\partial x^\mu}.
\label{eq:sr7-gradient}
\end{equation}
То, какое семейство называется верхним, а какое нижним, есть вопрос соглашения, подобно тому как соглашением являются названия «левый» и «правый». Однако после выбора этого соглашения его нужно последовательно соблюдать, потому что эти два семейства преобразуются по-разному.

\begin{classroomqa}
\audiencequestion{Когда индекс следует писать как верхний, а когда как нижний?}
\lecturerresponse{Координатам, смещениям и объектам, преобразующимся так же, как они, приписываются верхние индексы. Производным и объектам, преобразующимся так же, как производные, приписываются нижние индексы. Выбор того, какое семейство называть верхним, условен, но различие между этими двумя законами преобразования — нет. \lecturetimestamp{00:05:03}}
\end{classroomqa}

Для скалярной функции \(s(x)\) полный дифференциал имеет вид
\begin{equation}
ds
=
\frac{\partial s}{\partial x^\mu}\,dx^\mu
=
\partial_\mu s\,dx^\mu.
\label{eq:sr7-total-differential}
\end{equation}
В развернутой записи это
\begin{equation}
ds
=
\frac{\partial s}{\partial t}\,dt
+\frac{\partial s}{\partial x}\,dx
+\frac{\partial s}{\partial y}\,dy
+\frac{\partial s}{\partial z}\,dz.
\end{equation}
Компактное выражение говорит в точности то же самое. Это не новый закон; это способ увидеть скалярную структуру, не выписывая в каждой строке 4 члена.

\begin{figure}[htbp]
\centering
\includegraphics[width=0.90\textwidth]{lecture_07_figure_01.png}
\caption{Скалярный дифференциал, собранный из градиента с нижним индексом и смещения с верхним индексом.}
\lectureframe{00:08:20}
\end{figure}

Теперь соглашение Эйнштейна становится очевидным: по индексу, повторяющемуся один раз сверху и один раз снизу, ведётся суммирование. Поэтому свёртка в
\eqref{eq:sr7-total-differential} является скаляром. Существует более глубокая геометрическая интерпретация векторов и ковекторов, но это обсуждение отложено до общей теории относительности.

\subsection{Поднятие индекса}

У каждого 4-вектора есть формы с верхними и нижними индексами, связанные метрикой:
\begin{equation}
A_\mu=\eta_{\mu\nu}A^\nu.
\end{equation}
Используя \eqref{eq:sr7-metric},
\begin{equation}
A_0=-A^0,
\qquad
A_m=A^m,
\qquad m=1,2,3.
\label{eq:sr7-lowering}
\end{equation}
Тем самым поднятие или опускание индекса меняет знак временной компоненты и не затрагивает пространственные компоненты. То же правило применяется к производной:
\begin{equation}
\begin{aligned}
\partial_\mu
&=
\left(
\frac{\partial}{\partial t},
\frac{\partial}{\partial x},
\frac{\partial}{\partial y},
\frac{\partial}{\partial z}
\right),\\[0.5em]
\partial^\mu
&=
\left(
-\frac{\partial}{\partial t},
\frac{\partial}{\partial x},
\frac{\partial}{\partial y},
\frac{\partial}{\partial z}
\right).
\end{aligned}
\label{eq:sr7-raised-gradient}
\end{equation}

\begin{figure}[htbp]
\centering
\includegraphics[width=0.90\textwidth]{lecture_07_figure_02.png}
\caption{Операторы производной с нижним и верхним индексами; знак меняет только временная компонента.}
\lectureframe{00:12:30}
\end{figure}

\begin{classroomqa}
\audiencequestion{Всегда ли частная производная несёт нижний индекс, и можно ли его также поднять?}
\lecturerresponse{Оператор \(\partial/\partial x^\mu\) естественным образом есть \(\partial_\mu\), то есть ковектор. Применение метрики даёт \(\partial^\mu=\eta^{\mu\nu}\partial_\nu\); в плоском пространстве-времени это просто меняет знак временной компоненты. \lecturetimestamp{00:11:00}}
\end{classroomqa}

Инвариантный интервал — привычный пример:
\begin{equation}
x_\mu x^\mu
=
x^2+y^2+z^2-t^2.
\label{eq:sr7-interval}
\end{equation}
Свёртка одного верхнего и одного нижнего индекса даёт лоренцев скаляр. Именно из-за этого простого правила распознавания мы так тщательно ведём обозначения.

\section{Четыре принципа построения теории}

Предположим, эксперимент обнаруживает новое явление и мы хотим сформулировать его теорию. Поиск организуют 4 принципа. Впоследствии их, возможно, придётся уточнить, но всякая успешная фундаментальная теория в настоящее время соответствует этой схеме.

\subsection{Принцип действия}

Первое требование — наличие действия. Для механики частицы,
\begin{equation}
I[x]
=
\int dt\,L(x,\dot x),
\label{eq:sr7-particle-action-general}
\end{equation}
а физическая траектория стационарна относительно вариаций, при которых её концы фиксированы:
\begin{equation}
\delta I=0.
\end{equation}
Получающиеся уравнения имеют вид
\begin{equation}
\frac{d}{dt}
\left(
\frac{\partial L}{\partial\dot x^m}
\right)
=
\frac{\partial L}{\partial x^m}.
\label{eq:sr7-EL}
\end{equation}

Принцип действия — не эстетическое украшение. Это структура, которая связывает симметрии с законами сохранения. Симметрия действия относительно сдвигов по времени даёт сохранение энергии; пространственные сдвиги дают сохранение импульса; родственные симметрии лежат в основе других сохраняющихся величин. Само по себе действие не является законом сохранения, но именно это та структура, благодаря которой симметрия превращается в закон сохранения.

\subsection{Локальность}

Второе требование — локальность. Если возмутить одну точку скрипичной струны, сначала отреагирует соседний участок. Затем отреагирует его сосед, и возмущение будет распространяться. В конечном счёте эффект может дойти до всей струны, но непосредственное кратковременное влияние локально.

Для частицы локальность кодируется лагранжианом, вычисляемым вдоль траектории. Он зависит от положения и скорости в один и тот же момент времени, а не от произвольных пар далёких моментов времени. Для поля \(\phi(x)\),
\begin{equation}
I[\phi]
=
\int d^4x\,
\mathcal L\!\left(\phi,\partial_\mu\phi\right).
\label{eq:sr7-local-field-action}
\end{equation}
Плотность в одном событии зависит от поля и его первых производных в этом же событии. Она не содержит прямого произведения поля здесь с полем в другой, никак не связанной удалённой точке.

\begin{figure}[htbp]
\centering
\includegraphics[width=0.90\textwidth]{lecture_07_figure_03.png}
\caption{Действия для частицы и поля, показанные под принципом локальности.}
\lectureframe{00:19:30}
\end{figure}

Вот почему локальные уравнения движения являются дифференциальными уравнениями. Производные сравнивают значения в бесконечно близких точках. Нелокальная альтернатива могла бы связывать ускорение сейчас с силой в совершенно другой момент времени или связывать скорости в 2 удалённых друг от друга момента с силой в третий. Такие законы математически мыслимы, но локальность их исключает.

\begin{classroomqa}
\audiencequestion{Противоречит ли квантовая запутанность этому принципу локальности?}
\lecturerresponse{Нет, не в том динамическом смысле, который здесь имеется в виду. Запутанность создаёт корреляции, но локальную операцию, выполняемую Алисой, нельзя использовать для изменения редуцированного состояния Боба или для мгновенной передачи ему управляемого сигнала. Аналогия с пяти- и десятицентовой монетами хорошо выявляет это различие: увидев одну монету, узнаёшь, какая коррелированная монета у другого человека, но на удалённую монету это не действует. Квантовые корреляции богаче этой классической аналогии, и всё же они не позволяют осуществлять сверхсветовую передачу сигналов. \lecturetimestamp{00:20:02}}
\end{classroomqa}

\subsection{Лоренц-инвариантность}

Третье требование состоит в том, чтобы законы имели одну и ту же форму во всех инерциальных системах отсчёта. Практическое правило построения просто: лагранжиан должен быть лоренцевым скаляром, чтобы действие было инвариантным. Свёртки верхнего--нижнего индекса, такие как \(\partial_\mu\phi\,\partial^\mu\phi\), служат строительными блоками.

В общей теории относительности это утверждение расширяется. Подынтегральное выражение является скалярной плотностью, а инвариантность расширяется от преобразований Лоренца к произвольным преобразованиям координат. Это уточнение не заменяет нынешний принцип; это его форма для искривлённого пространства-времени.

\begin{classroomqa}
\audiencequestion{Является ли лоренц-инвариантность лишь правилом специальной теории относительности, а не по-настоящему общим принципом?}
\lecturerresponse{Для негравитационной задачи в плоском пространстве-времени лоренц-инвариантность — верное утверждение. В искривлённом пространстве-времени скалярный лагранжиан заменяют соответствующей скалярной плотностью и требуют инвариантности при произвольных преобразованиях координат. Помимо этого уточнения сохраняется та же философия построения. \lecturetimestamp{00:24:16}}
\end{classroomqa}

\begin{classroomqa}
\audiencequestion{Каковы размерности лагранжиана частицы и плотности лагранжиана поля?}
\lecturerresponse{Лагранжиан частицы имеет размерность энергии. В \(I=\int dt\,d^3x\,\mathcal L\) пространственная мера даёт объём, так что \(\mathcal L\) имеет размерность энергии на единицу объёма. Интегрирование по пространству даёт обычный лагранжиан с размерностью энергии.
\lecturetimestamp{00:25:49}}
\end{classroomqa}

\subsection{Калибровочная инвариантность}

Четвёртое требование — калибровочная инвариантность. На этом этапе его намеренно оставляют загадочным. Один из слушателей высказывает полезную предварительную догадку: возможно, к лагранжиану можно добавлять некоторые вещи, не изменяя физики. Это близко к истине, но почти всю работу здесь выполняет слово «некоторые».

\begin{figure}[htbp]
\centering
\includegraphics[width=0.90\textwidth]{lecture_07_figure_04.png}
\caption{На доске собраны локальность, скалярное построение, инвариантность Лоренца и калибровочная инвариантность.}
\lectureframe{00:27:00}
\end{figure}

Мы откладываем определение, потому что оно уже содержится в расчёте для заряженной частицы. Повторение этого расчёта — не отступление в сторону. Оно покажет, почему мы вынуждены ввести векторный потенциал и почему после этого он содержит избыточность.

\section{Почему необходим векторный потенциал}

Для медленно движущейся заряженной частицы целевым уравнением служит знакомый закон силы Лоренца
\begin{equation}
m\mathbf a
=
e\left(
\mathbf E+\mathbf v\times\mathbf B
\right).
\label{eq:sr7-lorentz-target}
\end{equation}

\begin{figure}[htbp]
\centering
\includegraphics[width=0.90\textwidth]{lecture_07_figure_05.png}
\caption{Знакомый закон силы, записанный до повторного введения векторного потенциала.}
\lectureframe{00:29:08.870}
\end{figure}

В \eqref{eq:sr7-lorentz-target} векторный потенциал не появляется. Уравнения движения можно записать непосредственно через \(\mathbf E\) и \(\mathbf B\). Трудность носит более структурный характер: при предположениях локальности и инвариантности Лоренца, используемых в лекции, взаимодействие заряженной частицы нельзя записать в виде действия на мировой линии непосредственно только через эти поля, положение и скорость. Потенциал \(A_\mu\) нужен, чтобы построить локальную скалярную связь.

\begin{classroomqa}
\audiencequestion{Если в закон силы уже входят только \(\mathbf E\) и \(\mathbf B\), зачем вводить \(A_\mu\)?}
\lecturerresponse{Потому что уравнение движения — это ещё не принцип действия. В локальной лоренц-инвариантной схеме этой лекции действие заряженной частицы требует четырёхпотенциала. Одни только поля позволяют сформулировать силу, но именно потенциал даёт возможность вывести её из скалярного действия на мировой линии. \lecturetimestamp{00:29:06}}
\end{classroomqa}

\subsection{Действие на мировой линии}

Предлагаемое действие имеет свободную часть и часть взаимодействия:
\begin{equation}
I[x]
=
-m\int d\tau
-e\int A_\mu(x)\,dx^\mu.
\label{eq:sr7-worldline-action}
\end{equation}
Первый член суммирует собственновременную длину каждого бесконечно малого участка траектории. По построению это скаляр. Во втором члене \(A_\mu(x)\) берётся в событии, которое занимает частица, и свёртывается со смещением \(dx^\mu\). Получающийся элемент тоже является скаляром и локален вдоль мировой линии.

\begin{figure}[htbp]
\centering
\includegraphics[width=0.90\textwidth]{lecture_07_figure_06.png}
\caption{Две части действия заряженной частицы и локальная свёртка \(A_\mu(x)dx^\mu\).}
\lectureframe{00:31:20}
\end{figure}

\subsection{Переписывание действия через обычное время}

Чтобы воспользоваться знакомыми уравнениями Эйлера--Лагранжа, разделим каждое приращение вдоль мировой линии на \(dt\):
\begin{equation}
\frac{d\tau}{dt}
=
\sqrt{1-\dot{\mathbf x}^{\,2}},
\qquad
\dot{\mathbf x}^{\,2}
=
\dot x^2+\dot y^2+\dot z^2.
\end{equation}
Четыре компоненты \(dx^\mu/dt\) имеют вид
\begin{equation}
\frac{dx^\mu}{dt}
=
\left(1,\dot x,\dot y,\dot z\right).
\end{equation}
Следовательно,
\begin{align}
I
&=
\int dt\,
\left[
-m\sqrt{1-\dot{\mathbf x}^{\,2}}
-eA_\mu\frac{dx^\mu}{dt}
\right]
\notag\\
&=
\int dt\,
\left[
-m\sqrt{1-\dot{\mathbf x}^{\,2}}
-e\left(A_0+A_m\dot x^m\right)
\right].
\label{eq:sr7-action-in-t}
\end{align}
Поэтому лагранжиан в обычном времени равен
\begin{equation}
\begin{aligned}
L(x,\dot x,t)
&=
-m\sqrt{1-\dot{\mathbf x}^{\,2}}
\\
&\quad -e\left[
A_0(x,t)+A_m(x,t)\dot x^m
\right].
\end{aligned}
\label{eq:sr7-lagrangian}
\end{equation}

\begin{figure}[htbp]
\centering
\includegraphics[width=0.90\textwidth]{lecture_07_figure_07.png}
\caption{Действие на мировой линии, переписанное как интеграл по обычному времени, и лагранжиан, считанный с его подынтегрального выражения.}
\lectureframe{00:35:30}
\end{figure}

На этом этапе творческая часть закончена. Как только лагранжиан выбран, всё остальное — механика. Эти выкладки могут быть длинными и утомительными, но эта утомительность не бессодержательна: именно она выявляет электрическую, магнитную и релятивистскую структуры. Лектор шутит, что открыть принцип — приятная часть дела, а выводить все его следствия — то, что поручают аспирантам.

\begin{classroomqa}
\audiencequestion{Можно ли в этом вычислении поднимать или опускать пространственный индекс у \(A_m\) или \(\dot x^m\)?}
\lecturerresponse{Да. Пространственная метрика единична, так что поднятие или опускание латинского индекса не меняет значения его компоненты. Осторожность со знаком требуется именно для временного индекса. \lecturetimestamp{00:34:10}}
\end{classroomqa}

\section{Вычисление по уравнениям Эйлера--Лагранжа}

\subsection{Канонический импульс и исправление по ходу дела}

Выберем одну пространственную компоненту \(x^m\). Дифференцирование свободного члена в \eqref{eq:sr7-lagrangian} даёт
\begin{equation}
\frac{m\dot x_m}{\sqrt{1-\dot{\mathbf x}^{\,2}}}
=
m u_m,
\qquad
u^\mu\equiv\frac{dx^\mu}{d\tau}.
\end{equation}
На доске первая попытка на этом останавливается. Замечание из аудитории указывает на пропуск: член взаимодействия тоже зависит от \(\dot x^m\). Полная производная равна
\begin{equation}
\frac{\partial L}{\partial\dot x^m}
=
\frac{m\dot x_m}{\sqrt{1-\dot{\mathbf x}^{\,2}}}
-eA_m
=
m u_m-eA_m.
\label{eq:sr7-canonical-momentum}
\end{equation}

\begin{classroomqa}
\audiencequestion{Не отсутствует ли в первом выражении для \(\partial L/\partial\dot x^m\) вклад от электромагнитного члена?}
\lecturerresponse{Да. Дифференцирование
\(-eA_n\dot x^n\) по \(\dot x^m\) даёт вклад
\(-eA_m\). Исправленный канонический импульс равен
\(m u_m-eA_m\), а не только механическому импульсу.
\lecturetimestamp{00:37:24}}
\end{classroomqa}

\begin{figure}[htbp]
\centering
\includegraphics[width=0.90\textwidth]{lecture_07_figure_08.png}
\caption{Исправленный канонический импульс с механическим и потенциальным членами.}
\lectureframe{00:38:30}
\end{figure}

\subsection{Дифференцирование вдоль движущейся траектории}

Левая часть уравнения Эйлера--Лагранжа имеет вид
\begin{equation}
\frac{d}{dt}
\left(
\frac{\partial L}{\partial\dot x^m}
\right)
=
m\frac{du_m}{dt}
-e\frac{dA_m(x(t),t)}{dt}.
\end{equation}
Потенциал меняется по двум причинам: он может явно зависеть от времени, и частица перемещается по области, где он меняется в пространстве. Правило цепочки даёт
\begin{equation}
\frac{dA_m}{dt}
=
\partial_0A_m
+\dot x^n\partial_nA_m.
\label{eq:sr7-chain-rule}
\end{equation}
В развёрнутой записи второй член означает
\begin{equation}
\dot x\,\frac{\partial A_m}{\partial x}
+\dot y\,\frac{\partial A_m}{\partial y}
+\dot z\,\frac{\partial A_m}{\partial z}.
\end{equation}

\begin{figure}[htbp]
\centering
\includegraphics[width=0.90\textwidth]{lecture_07_figure_09.png}
\caption{Явная и конвективная части производной \(A_m\) по времени.}
\lectureframe{00:42:10}
\end{figure}

В правой части,
\begin{equation}
\frac{\partial L}{\partial x^m}
=
-e\partial_mA_0
-e\dot x^n\partial_mA_n.
\label{eq:sr7-spatial-derivative-L}
\end{equation}
Свободный член явно от положения не зависит. В \eqref{eq:sr7-spatial-derivative-L} индекс \(m\) — свободный индекс, выбирающий одно уравнение; \(n\) — повторяющийся немой индекс, по которому ведётся суммирование. Переименование этого немого индекса ничего не меняет. Именно поэтому лектор, прежде чем объединить члены, останавливается на буквах.

Приравнивая обе части и перегруппировывая члены, получаем
\begin{equation}
m\frac{du_m}{dt}
=
e\left(
\partial_0A_m-\partial_mA_0
\right)
+e\dot x^n
\left(
\partial_nA_m-\partial_mA_n
\right).
\label{eq:sr7-grouped-force}
\end{equation}

\begin{figure}[htbp]
\centering
\includegraphics[width=0.90\textwidth]{lecture_07_figure_10.png}
\caption{Сила, разложенная на скобку, не зависящую от скорости, и антисимметричную скобку, линейную по скорости.}
\lectureframe{00:49:30}
\end{figure}

\subsection{Распознавание электрического и магнитного полей}

Первая скобка не содержит скорости частицы. В соглашении, принятом для этого действия, определим
\begin{equation}
E_m
\equiv
\partial_0A_m-\partial_mA_0.
\label{eq:sr7-E-from-A}
\end{equation}
Она играет в точности ту операционную роль, которую играет электрическое поле. Вторая скобка антисимметрична по своим пространственным индексам:
\begin{equation}
C_{mn}
\equiv
\partial_nA_m-\partial_mA_n.
\label{eq:sr7-C}
\end{equation}
Её три независимые компоненты образуют ротор пространственного потенциала; ориентация выбрана в соответствии с соглашением лекции так, что
\begin{equation}
\dot x^n C_{mn}
=
\left(\mathbf v\times\mathbf B\right)_m.
\label{eq:sr7-magnetic-identification}
\end{equation}
Подробная покомпонентная проверка оставлена на доске в качестве упражнения. Структурное отождествление уже видно: антисимметричная пространственная производная вектора имеет три компоненты — ровно столько, сколько имеет \(\mathbf B\), — и она умножается на скорость.

\begin{figure}[htbp]
\centering
\includegraphics[width=0.90\textwidth]{lecture_07_figure_11.png}
\caption{Две скобки с производными, отождествлённые с электрической и магнитной частями силы Лоренца.}
\lectureframe{00:52:00}
\end{figure}

Тем самым уравнение \eqref{eq:sr7-grouped-force} восстановило силу Лоренца из действия с векторным потенциалом. Но оно сделало и нечто более важное: выявило антисимметризованные производные, которые станут тензором поля.

\section{От обычного времени к четырёхвекторному уравнению}

\subsection{Переход к собственному времени}

Левая часть \eqref{eq:sr7-grouped-force} — это не ньютоновское ускорение \(d\mathbf v/dt\). Это производная компоненты четырёхскорости по обычному времени:
\begin{equation}
\frac{du^m}{dt}
=
\frac{du^m}{d\tau}\frac{d\tau}{dt}.
\end{equation}
Умножим уравнение на \(dt/d\tau\). Тогда
\begin{equation}
\dot x^n\frac{dt}{d\tau}
=
\frac{dx^n}{dt}\frac{dt}{d\tau}
=
u^n,
\qquad
\frac{dt}{d\tau}=u^0.
\end{equation}
Пространственное уравнение принимает вид
\begin{align}
m\frac{du^m}{d\tau}
&=
e\left(
\partial_0A_m-\partial_mA_0
\right)u^0
\notag\\
&\quad
+e\left(
\partial_nA_m-\partial_mA_n
\right)u^n.
\label{eq:sr7-proper-spatial}
\end{align}

\begin{figure}[htbp]
\centering
\includegraphics[width=0.90\textwidth]{lecture_07_figure_12.png}
\caption{Множители \(dt/d\tau\) превращают обычные скорости в компоненты четырёхскорости.}
\lectureframe{00:56:30}
\end{figure}

Левая часть есть собственное ускорение:
\begin{equation}
\frac{du^\mu}{d\tau}
=
\frac{d^2x^\mu}{d\tau^2}.
\end{equation}
Это четырёхвектор. Следовательно, правая часть должна быть построена из тензора второго ранга, действующего на \(u^\nu\).

\subsection{Тензор поля и закон силы}

Примем соглашение
\begin{equation}
F_{\mu\nu}
\equiv
\partial_\mu A_\nu-\partial_\nu A_\mu.
\label{eq:sr7-F-definition}
\end{equation}
Он антисимметричен:
\begin{equation}
F_{\mu\nu}=-F_{\nu\mu}.
\end{equation}
С учетом знака взаимодействия в \eqref{eq:sr7-worldline-action} компактное уравнение имеет вид
\begin{equation}
\boxed{
m\frac{du^\mu}{d\tau}
=
-eF^\mu{}_\nu u^\nu
}.
\label{eq:sr7-covariant-force}
\end{equation}
Эквивалентно, до введения обозначения \(F_{\mu\nu}\),
\begin{equation}
m\frac{du^\mu}{d\tau}
=
e\left(
\partial_\nu A^\mu-\partial^\mu A_\nu
\right)u^\nu.
\end{equation}
Различные соглашения о метрике, заряде и потенциале перераспределяют общие знаки. Инвариантное содержание составляют антисимметризованная производная потенциала и ее свертка с четырёхскоростью.

\begin{figure}[htbp]
\centering
\includegraphics[width=0.90\textwidth]{lecture_07_figure_13.png}
\caption{Уравнение по собственному времени, записанное в виде тензора поля второго ранга, действующего на четырёхскорость.}
\lectureframe{00:58:25}
\end{figure}

\begin{classroomqa}
\audiencequestion{При выводе получились три пространственных уравнения. Почему допустимо добавить четвертое?}
\lecturerresponse{Скалярное действие гарантирует лоренц-ковариантный результат. Когда разность между двумя сторонами распознана как четырёхвектор, преобразования Лоренца смешивают его пространственные и временную компоненты. Уравнение для четырёхвектора не может согласованно иметь только три компоненты; временная компонента сопровождает пространственные. \lecturetimestamp{00:59:40}}
\end{classroomqa}

\begin{classroomqa}
\audiencequestion{Является ли сам по себе фиксированный набор чисел, такой как \((1,2,3,4)\), лоренц-инвариантным четырёхвектором?}
\lecturerresponse{Нет. Такие числа могут быть компонентами вектора в одной выбранной системе отсчета, но компоненты преобразуются. Даже в двух измерениях вектор с компонентами \((1,2)\) можно описать в повернутой системе отсчета как \((0,\sqrt5)\). Геометрический вектор сохраняется; его численный список компонент — нет. \lecturetimestamp{01:00:28}}
\end{classroomqa}

\begin{figure}[htbp]
\centering
\includegraphics[width=0.90\textwidth]{lecture_07_figure_14.png}
\caption{Двухкомпонентный вектор, переписанный в осях, направленных вдоль самого вектора.}
\lectureframe{01:02:00}
\end{figure}

\begin{classroomqa}
\audiencequestion{Если цель состояла в том, чтобы заменить обычное время собственным, почему все еще появляется \(\partial_0A_m=\partial A_m/\partial t\)?}
\lecturerresponse{Цель не состояла в том, чтобы стереть всякую букву \(t\). Она состояла в том, чтобы записать динамику в явной четырёхвекторной нотации. Оператор \(\partial_0\) является временной компонентой ковариантного четырёхградиента, так же как \(\partial_m\) являются его пространственными компонентами. Ему место внутри четырёхмерного тензора. \lecturetimestamp{01:05:10}}
\end{classroomqa}

Лекция использует этот момент, чтобы заново сформулировать локальность. Локальное уравнение связывает ускорение в некотором событии с полем и скоростью в этом событии. Оно не связывает ускорение сейчас с силой в удаленный момент времени.

\begin{figure}[htbp]
\centering
\includegraphics[width=0.90\textwidth]{lecture_07_figure_15.png}
\caption{Ковариантное уравнение силы, использованное для нового формулирования локальности как дифференциальной динамики.}
\lectureframe{01:05:50}
\end{figure}

\subsection{Явный вид электромагнитных компонент}

Для тензора со всеми верхними индексами соглашение, использованное на доске, дает
\begin{equation}
F^{\mu\nu}
=
\begin{pmatrix}
0&-E_x&-E_y&-E_z\\
E_x&0&B_z&-B_y\\
E_y&-B_z&0&B_x\\
E_z&B_y&-B_x&0
\end{pmatrix}.
\label{eq:sr7-F-matrix}
\end{equation}
Индекс строки — \(\mu\), индекс столбца — \(\nu\). Смешанные временные--пространственные элементы содержат электрическое поле; пространственный--пространственный блок содержит магнитное поле. Антисимметрия обнуляет диагональ и сводит 16 выписанных элементов к 6 независимым компонентам.

\begin{figure}[htbp]
\centering
\includegraphics[width=0.90\textwidth]{lecture_07_figure_16.png}
\caption{Полный электромагнитный тензор: электрические элементы стоят во временных--пространственных позициях, а магнитные — в пространственном блоке.}
\lectureframe{01:12:10}
\end{figure}

\begin{classroomqa}
\audiencequestion{Существует ли стандартное правило, определяющее, какой индекс тензора обозначает строки, а какой — столбцы?}
\lecturerresponse{Это соглашение о способе записи. Здесь первый индекс \(\mu\) идет по вертикали, а второй индекс \(\nu\) — по горизонтали. После объявления этого соглашения каждую компоненту можно читать однозначно. Поднятие временного индекса вносит знак, но не меняет соглашение строка--столбец. \lecturetimestamp{01:09:01}}
\end{classroomqa}

\begin{classroomqa}
\audiencequestion{Какие расположения индексов сохраняют антисимметрию?}
\lecturerresponse{Тензоры \(F_{\mu\nu}\) и \(F^{\mu\nu}\) антисимметричны. Смешанный массив \(F^\mu{}_\nu\) не является антисимметричным при обычном транспонировании матрицы: поднятие одного временного индекса меняет электрические ячейки, делая временную--пространственную пару симметричной, тогда как чисто пространственный магнитный блок остается антисимметричным. Лектор приходит к этому выводу после того, как ненадолго проверяет на доске противоположную возможность и отвергает ее.
\lecturetimestamp{01:13:18}}
\end{classroomqa}

\section{Что считается инвариантностью?}

Только теперь мы возвращаемся к четвертому принципу. Инвариантность — это изменение, оставляющее уравнения движения неизменными. Само действие может остаться в точности неизменным, а может измениться на граничный член, который не влияет на вариационные уравнения.

\subsection{Постоянные сдвиги безмассового скалярного поля}

Рассмотрим уравнение безмассового скалярного поля в одном пространственном измерении,
\begin{equation}
\frac{\partial^2\phi}{\partial t^2}
-\frac{\partial^2\phi}{\partial x^2}
=0,
\label{eq:sr7-wave}
\end{equation}
получаемое из
\begin{equation}
\mathcal L_0
=
-\frac12\partial_\mu\phi\,\partial^\mu\phi.
\label{eq:sr7-massless-L}
\end{equation}
Если \(\phi\) является решением, то решением является и
\begin{equation}
\phi\longrightarrow\phi+c,
\label{eq:sr7-constant-shift}
\end{equation}
где \(c\) постоянно. Все производные от \(c\) обращаются в нуль, поэтому ни уравнение, ни лагранжиан, содержащий только производные, не меняются.

\begin{figure}[htbp]
\centering
\includegraphics[width=0.90\textwidth]{lecture_07_figure_17.png}
\caption{Безмассовое волновое уравнение, его лагранжиан, зависящий только от производных, и постоянный сдвиг \(\phi\mapsto\phi+c\).}
\lectureframe{01:17:10}
\end{figure}

Теперь добавим массовый член:
\begin{equation}
\frac{\partial^2\phi}{\partial t^2}
-\frac{\partial^2\phi}{\partial x^2}
+\mu^2\phi
=0,
\end{equation}
\begin{equation}
\mathcal L_\mu
=
-\frac12\partial_\mu\phi\,\partial^\mu\phi
-\frac{\mu^2}{2}\phi^2.
\label{eq:sr7-massive-L}
\end{equation}
Часть с производными остается неизменной при
\(\phi\mapsto\phi+c\), но \(\phi^2\) — нет. Обе теории допустимы; одна обладает инвариантностью относительно сдвига, а другая — нет.

\begin{figure}[htbp]
\centering
\includegraphics[width=0.90\textwidth]{lecture_07_figure_18.png}
\caption{Массовый член, нарушающий инвариантность относительно постоянного сдвига.}
\lectureframe{01:18:50}
\end{figure}

\subsection{Студент предвосхищает граничный член}

Один из слушателей предлагает более строгую проверку:
\begin{equation}
\phi(x,t)\longrightarrow\phi(x,t)+x.
\label{eq:sr7-linear-shift}
\end{equation}
Вторые производные от \(x\) обращаются в нуль, поэтому безмассовое волновое уравнение по-прежнему выполняется. Действие требует более внимательного рассмотрения. Поскольку \(x\) и \(t\) являются независимыми переменными,
\begin{equation}
\frac{\partial x}{\partial t}=0,
\qquad
\frac{\partial x}{\partial x}=1.
\end{equation}
Пространственная часть изменяется согласно
\begin{align}
-\frac12
\left[
\partial_x(\phi+x)
\right]^2
&=
-\frac12
\left(
\partial_x\phi+1
\right)^2
\notag\\
&=
-\frac12(\partial_x\phi)^2
-\partial_x\phi
-\frac12.
\label{eq:sr7-linear-shift-L}
\end{align}
Новые слагаемые — это константа и полная производная. Константа не изменяет уравнение Эйлера--Лагранжа. Интегрирование \(\partial_x\phi\) дает только пространственные граничные значения, так что при обычных фиксированных граничных условиях оно также не изменяет уравнение движения.

\begin{figure}[htbp]
\centering
\includegraphics[width=0.90\textwidth]{lecture_07_figure_19.png}
\caption{Предложенный сдвиг \(\phi\mapsto\phi+x\) и квадрат производной, из которого видны константа и полная производная.}
\lectureframe{01:22:40}
\end{figure}

\begin{classroomqa}
\audiencequestion{Сохраняет ли \(\phi\mapsto\phi+x\) безмассовое уравнение поля, хотя и изменяет плотность лагранжиана?}
\lecturerresponse{Да. Вторые производные от \(x\) обращаются в нуль, так что уравнение не меняется. Лагранжиан изменяется только на константу и полную пространственную производную. Они изменяют численное значение действия лишь на фиксированную константу по области интегрирования и на граничные значения, но не уравнение Эйлера--Лагранжа при фиксированных граничных условиях. \lecturetimestamp{01:20:00}}
\end{classroomqa}

\begin{classroomqa}
\audiencequestion{Почему \(\partial x/\partial t\) равно нулю?}
\lecturerresponse{Потому что \(x\) и \(t\) — независимые координаты в описании поля. Частная производная по \(t\) удерживает \(x\) фиксированным. \lecturetimestamp{01:22:22}}
\end{classroomqa}

Студент пришел к приему, который нужен для электромагнетизма: полная производная может изменить записанное действие, не меняя его стационарных траекторий.

\section{Калибровочная инвариантность заряженной частицы}

\subsection{Добавление четырёхградиента}

Вернемся к взаимодействию
\begin{equation}
I_{\mathrm{int}}
=
-e\int A_\mu\,dx^\mu.
\end{equation}
Пусть \(\Lambda(x)\) — любой достаточно гладкий скаляр, и изменим потенциал по правилу
\begin{equation}
A_\mu
\longrightarrow
A'_\mu
=
A_\mu+\partial_\mu\Lambda.
\label{eq:sr7-gauge-transformation}
\end{equation}
Градиент \(\partial_\mu\Lambda\) имеет ту же ковариантную индексную структуру, что и \(A_\mu\), так что сумма определена корректно.

\begin{figure}[htbp]
\centering
\includegraphics[width=0.90\textwidth]{lecture_07_figure_20.png}
\caption{Связь с мировой линией и преобразование \(A_\mu\mapsto A_\mu+\partial_\mu\Lambda\).}
\lectureframe{01:24:50}
\end{figure}

Изменение действия равно
\begin{align}
\Delta I_{\mathrm{int}}
&=
-e\int \partial_\mu\Lambda\,dx^\mu
\notag\\
&=
-e\int d\Lambda
\notag\\
&=
-e\left[
\Lambda(x_2)-\Lambda(x_1)
\right].
\label{eq:sr7-gauge-boundary}
\end{align}
На каждом бесконечно малом участке мировой линии \(\partial_\mu\Lambda\,dx^\mu\) есть в точности приращение
\(d\Lambda\). Суммирование приращений телескопически сводится к разности конечного и начального значений на концах.

\begin{figure}[htbp]
\centering
\includegraphics[width=0.90\textwidth]{lecture_07_figure_21.png}
\caption{Вклад градиента, проинтегрированный вдоль пути, дает разность на концах \(\Lambda(x_2)-\Lambda(x_1)\).}
\lectureframe{01:28:10}
\end{figure}

Вариационный принцип сравнивает близкие траектории при фиксированных \(x_1\) и
\(x_2\). Поэтому дополнительный член в \eqref{eq:sr7-gauge-boundary} одинаков для любой допустимой вариации:
\begin{equation}
\delta\Delta I_{\mathrm{int}}=0.
\end{equation}
Стационарная траектория не меняется. Это калибровочная инвариантность как утверждение о конечных точках.

\begin{classroomqa}
\audiencequestion{Может ли \(\Lambda\) быть любой скалярной функцией?}
\lecturerresponse{Да, при соблюдении предположений о гладкости и граничном поведении, необходимых для производных и действия. Преобразование локально, потому что параметр является произвольной функцией пространства-времени, а не просто одной глобальной константой. \lecturetimestamp{01:28:40}}
\end{classroomqa}

Добавление градиента к \(A_\mu\) называется калибровочным преобразованием. Слово gauge — историческое наследие; лекция с юмором, но вполне твердо, предостерегает от того, чтобы произносить его так, как будто это французское слово. Название менее важно, чем сама операция.

\subsection{Прямое доказательство с помощью тензора поля}

Уравнения движения зависят от \(F_{\mu\nu}\), а не от одного только потенциала. При преобразовании \eqref{eq:sr7-gauge-transformation},
\begin{align}
F'_{\mu\nu}
&=
\partial_\mu
\left(
A_\nu+\partial_\nu\Lambda
\right)
-
\partial_\nu
\left(
A_\mu+\partial_\mu\Lambda
\right)
\notag\\
&=
F_{\mu\nu}
+
\partial_\mu\partial_\nu\Lambda
-
\partial_\nu\partial_\mu\Lambda
\notag\\
&=
F_{\mu\nu}.
\label{eq:sr7-F-gauge-invariant}
\end{align}
Для гладкого скаляра в обычном плоском пространстве-времени частные производные коммутируют.
Два дополнительных члена сокращаются. Следовательно, тензор поля, содержащиеся в нём электрическое и
магнитное поля и закон силы — всё это калибровочно-инвариантно.

\begin{figure}[htbp]
\centering
\includegraphics[width=0.90\textwidth]{lecture_07_figure_22.png}
\caption{Две вторые производные, возникающие при калибровочном преобразовании, сокращаются, поскольку частные производные коммутируют.}
\lectureframe{01:33:50}
\end{figure}

Два доказательства замыкают один и тот же круг, идя с противоположных сторон. Действие
изменяется лишь в фиксированных конечных точках, тогда как уравнения зависят только от
неизменившегося тензора поля.

\section{Что означает избыточность}

\subsection{Зачем сохранять объект, который не определён однозначно?}

Естественное возражение возникает сразу. Если \(A_\mu\) не определён однозначно, зачем мы
его ввели? Потому что он необходим для локального лоренц-инвариантного
действия, хотя значения его компонент содержат избыточную информацию.
Одно и то же электромагнитное поле описывается многими потенциалами:
\begin{equation}
A_\mu
\sim
A_\mu+\partial_\mu\Lambda.
\end{equation}
Физически значим именно класс эквивалентности; конкретный представитель — это выбор
описания.

\begin{classroomqa}
\audiencequestion{Есть ли практическая польза от смены калибровки?}
\lecturerresponse{Да. Подходящая \(\Lambda\) может упростить одну из компонент
потенциала или сделать некоторую особенность задачи наглядной. Локально и
при подходящих граничных условиях часто можно выбрать калибровку, в которой
одна выбранная компонента, в том числе \(A_0\), обращается в нуль. Другая калибровка может
яснее выявить иную особенность. \lecturetimestamp{01:35:43}}
\end{classroomqa}

\begin{classroomqa}
\audiencequestion{Аналогична ли калибровочная инвариантность описанию одного и того же эксперимента из двух лоренцевых систем отсчёта?}
\lecturerresponse{Нет. Преобразования Лоренца сопоставляют физически
осмысленные системы отсчёта, связанные с разными наблюдателями. Калибровочные
преобразования — это избыточности внутри одного описания: разные
потенциалы кодируют одно и то же поле. Математически и то и другое является инвариантностью, но
их интерпретации различны. \lecturetimestamp{01:37:07}}
\end{classroomqa}

Есть и структурное различие. Вращение использует фиксированный глобальный
угол; обычно не вращают оси на несвязанные друг с другом углы в каждой
точке. Калибровочное преобразование параметризуется всей функцией
\(\Lambda(x)\). Именно эта зависящая от положения свобода и делает такую
инвариантность локальной калибровочной.

\begin{classroomqa}
\audiencequestion{Является ли добавление градиента общим определением калибровочной инвариантности?}
\lecturerresponse{Для абелевой электромагнитной теории именно так выглядит
преобразование. В теориях Yang--Mills есть более сложная неабелева
версия, в которой потенциал несёт дополнительную внутреннюю структуру, но
основная идея сохраняется: разные локальные описания описывают одну и ту же
физику. \lecturetimestamp{01:39:27}}
\end{classroomqa}

\begin{classroomqa}
\audiencequestion{Является ли калибровочное преобразование всего лишь математическим приёмом, позволяющим выделять разные стороны задачи?}
\lecturerresponse{Это полезный вычислительный и интерпретационный инструмент, но
его роль ограничена: оно может упростить уравнения или выявить
некоторое свойство, однако не может изменить никакого окончательного
калибровочно-инвариантного предсказания.
\lecturetimestamp{01:40:22}}
\end{classroomqa}

\subsection{Другие поля и другие потенциалы}

Калибровочные потенциалы в других теориях тоже обычно обозначают через \(A\) и
снабжают пространственно-временным индексом \(\mu\). Их отличает дополнительная
внутренняя структура, связанная с другими симметриями. В лекции обсуждение
этой структуры откладывается, вместо того чтобы предлагать преждевременную формулу.

\begin{classroomqa}
\audiencequestion{Как выглядят векторные потенциалы в других теориях поля?}
\lecturerresponse{Они по-прежнему ведут себя как пространственно-временные векторы и обычно
записываются как \(A_\mu\), но несут дополнительные метки и алгебраическую структуру,
связанную с внутренними симметриями. С этого начинается теория Yang--Mills,
и её обсуждение намеренно отложено. \lecturetimestamp{01:41:02}}
\end{classroomqa}

\begin{classroomqa}
\audiencequestion{Неужели действительно невозможно записать локальное лоренц-инвариантное действие заряженной частицы напрямую через \(\mathbf E\) и \(\mathbf B\)?}
\lecturerresponse{В рамках сформулированных здесь предпосылок — да. В лекции
это утверждается как теорема и не доказывается. Сохранение принципа
действия, локальности и лоренц-инвариантности вынуждает ввести
потенциал, после чего калибровочная избыточность возникает автоматически.
\lecturetimestamp{01:42:18}}
\end{classroomqa}

\subsection{Что можно измерить?}

Локальное значение \(A_\mu\) само по себе не является калибровочно-инвариантным. Если бы физическое
измерение однозначно определяло это значение, то его изменение на
\(\partial_\mu\Lambda\) меняло бы физический результат, вопреки только что
доказанной инвариантности. Измеримые предсказания должны строиться из
калибровочно-инвариантных величин, таких как \(F_{\mu\nu}\) в настоящем классическом
обсуждении.

\begin{classroomqa}
\audiencequestion{Почему природа так упорно выбирает калибровочные теории?}
\lecturerresponse{На классическом уровне лекция не предлагает окончательной причины.
Квантовая механика даёт более глубокую перспективу, но даже там общий
ответ не исчерпан. Уже сейчас очевидны необычная красота,
симметрия и экономность теорий, организованных калибровочной избыточностью.
\lecturetimestamp{01:43:36}}
\end{classroomqa}

\begin{classroomqa}
\audiencequestion{Можно ли измерить сам векторный потенциал?}
\lecturerresponse{Не как калибровочно-зависимое локальное значение компоненты.
Измерение — физический процесс, и всякое физическое предсказание должно быть
нечувствительно к \(A_\mu\mapsto A_\mu+\partial_\mu\Lambda\). Измеримы калибровочно-инвариантные свойства потенциала и поля.
\lecturetimestamp{01:44:18}}
\end{classroomqa}

\begin{classroomqa}
\audiencequestion{Что происходит, если \(A_\mu\) изменить на нечто, не являющееся градиентом?}
\lecturerresponse{Вообще говоря, \(F_{\mu\nu}\) изменяется, поэтому такое изменение
обнаружимо. Напротив, чисто калибровочный потенциал
\(A_\mu=\partial_\mu\Lambda\) имеет
\(F_{\mu\nu}=0\) и не оказывает электромагнитной силы; тогда свободная частица
движется по прямой мировой линии. \lecturetimestamp{01:45:22}}
\end{classroomqa}

Поэтому калибровочная инвариантность — не трюк, от которого можно просто отказаться. Электродинамика,
Yang--Mills, Стандартная модель и гравитация все обладают калибровочными
избыточностями того или иного рода. Эта схема своеобразна: для формулировки
теории требуется вспомогательная величина, однако часть этой величины можно
изменять без изменения физики. Это своеобразие и есть
организующая сила.

\section{Логика калибровочной инвариантности}

Лекция пошла по намеренно непрямому пути. Сначала мы восстановили
язык верхних--нижних индексов, поскольку скалярные свёртки — это грамматика
лоренц-инвариантных действий. Затем мы сформулировали четыре принципа: действие,
локальность, лоренц-инвариантность и калибровочную инвариантность.
Четвёртый оставался неопределённым, пока заново выстраивалось вычисление для заряженной частицы.

Это вычисление началось с локального скалярного действия
\begin{equation}
I
=
-m\int d\tau
-e\int A_\mu dx^\mu,
\end{equation}
прошло через лагранжиан по обычному времени и все производные Эйлера--Лагранжа,
и восстановило электрическую и магнитную силу Лоренца. Затем собственное
время свело результат к виду
\begin{equation}
m\frac{du^\mu}{d\tau}
=
-eF^\mu{}_\nu u^\nu,
\qquad
F_{\mu\nu}
=
\partial_\mu A_\nu-\partial_\nu A_\mu.
\end{equation}

Примеры со скалярным полем подготовили последний шаг: полная производная может
изменить действие только через границы. Следовательно,
\begin{equation}
A_\mu\longrightarrow A_\mu+\partial_\mu\Lambda
\end{equation}
изменяет действие мировой линии лишь на фиксированные значения в конечных точках и оставляет
\(F_{\mu\nu}\) неизменным, потому что смешанные частные производные коммутируют. Поэтому
векторный потенциал необходим, но не определён однозначно. Калибровочная инвариантность —
это утверждение, что эта неоднозначность относится к описанию, а не к физике.
